\section{Integral and $\lambda$-ring controls on the bridges}\label{sec:root-findings}
The results in this section are deductions checked in this run. They are not claims of historical priority. We distinguish an identity of characters, a decomposition in a Grothendieck group, an actual direct-sum decomposition, and a dimension identity.

\begin{proposition}[G-01: The three meanings of the second operation]\label{G-01}
Let $E$ be a finite $G$-set. Write $b_2(E)$ for the $G$-set of two-element subsets and $s_2(E)$ for the $G$-set of two-element multisets. In the Burnside ring one has
\[
 [s_2(E)]=[b_2(E)]+[E],\qquad [E]^2=[E]+[\operatorname{Inj}(\{1,2\},E)].
\]
The identity $[E]^2=2[b_2(E)]+[E]$ does not hold in general. If $V$ is a representation over a field of characteristic zero, the correct special $\lambda$-ring identities are
\[
 [\operatorname{Sym}^2V]=[V]^2-[\wedge^2V],\quad
 [\operatorname{Sym}^2V]-[\wedge^2V]=\psi^2[V],\quad
 [V]^2=2[\wedge^2V]+\psi^2[V].
\]
Consequently the expression $2\lambda^2(E)+E=E^2$ in cote 133, scan 019, line 53 is valid after taking cardinalities, but is not a universal equivariant or special-$\lambda$-ring identity. The supplied study's assertion that the two equivariant descriptions agree needs this correction.
\end{proposition}
\begin{proof}
Partition a multiset of cardinality two into the diagonal multisets $\{e,e\}$ and the subsets $\{e,e'\}$ with $e\ne e'$. This partition is $G$-equivariant. Similarly partition ordered pairs into diagonal and off-diagonal pairs. For a counterexample take $G=C_2$ and its regular two-element set $E$. Then $E\times E$ is a disjoint union of two free $C_2$-orbits, whereas $b_2(E)$ is one fixed point. Thus the asserted equality would read $2[E]=[E]+2[\mathrm{pt}]$. Applying the fixed-point homomorphism for the whole group gives $0=2$, a contradiction.
For representations, first prove the universal identity on $\mathrm{GL}_r$ and then pull it back along the representation of $G$. Write the character of its defining representation on the diagonal torus as $x_1+\cdots+x_r$. The symmetric and exterior squares have characters $\sum_{i\leq j}x_ix_j$ and $\sum_{i<j}x_ix_j$. Their sum is $(\sum_i x_i)^2$, and their difference is $\sum_i x_i^2=\psi^2[V]$. Character injectivity in characteristic zero proves the formulas. In the same $C_2$ example the permutation representation is $1\oplus\epsilon$, and $\wedge^2V=\epsilon$, whereas linearization of $b_2(E)$ is $1$. This explicitly shows why the subset operation cannot be silently identified with exterior power. Taking dimensions sends $\psi^2[V]$ to $\dim V$ and recovers the numerical identity. None of this determines the ink on a scan which has not been supplied.
\end{proof}

\begin{proposition}[G-02: The terminal exterior-power variable is necessary]\label{G-02}
For $i\geq2$, a universal formula for $\lambda^i(xy)$ in special $\lambda$-rings cannot in general omit $\lambda^i(x)$ or $\lambda^i(y)$. Thus the terminal index in cote 133, scan 006, lines 25--26 must be $i$ if those lines state the universal product formula.
\end{proposition}
\begin{proof}
Set $y=1$, for which $\lambda^1(1)=1$ and $\lambda^j(1)=0$ for $j>1$. A formula using only $x,\lambda^2(x),\ldots,\lambda^{i-1}(x)$ would express $\lambda^i(x)$ in the subring generated by these earlier operations. Evaluate at the defining representation of $\mathrm{GL}_i$ and restrict to its diagonal torus. This would express the elementary symmetric polynomial $e_i$ in $\mathbb Z[e_1,\ldots,e_{i-1}]$, contrary to the algebraic independence of $e_1,\ldots,e_i$ (equivalently the fundamental theorem of symmetric polynomials). Interchange $x$ and $y$ for the other endpoint. The argument establishes the required mathematical reading, not a new visual transcription.
\end{proof}

\begin{proposition}[G-03: Finite $\gamma$-data do not determine multiplicities]\label{G-03}
For every integer $D\geq0$ there are two honest representations $A_D,B_D$ of the split group scheme $\mathbb G_{m,\mathbb Z}$, each of rank $2^D$, which have the same class modulo $F^{D+1}_\gamma$ but different weight multiplicities. Over $\mathbb Q$ their Chern characters agree in degrees $0,\ldots,D$.
\end{proposition}
\begin{proof}
Let $L$ denote the standard character and put $r=D+1$. Define
\[
 A_D=\bigoplus_{\substack{0\leq j\leq r\\j\text{ even}}}L^{\otimes j\,\oplus\binom rj},\qquad
 B_D=\bigoplus_{\substack{0\leq j\leq r\\j\text{ odd}}}L^{\otimes j\,\oplus\binom rj}.
\]
The two binomial parity sums are $2^{r-1}$. In the representation ring $\mathbb Z[L,L^{-1}]$ the difference of the classes is $(1-L)^r$, a product of $r$ rank-zero elements and hence an element of $F^r_\gamma$. The constant weight has multiplicity one in $A_D$ and zero in $B_D$, so the characters differ. With $c=c_1(L)$, the difference of the Chern characters is $(1-e^c)^r$, whose lowest possible degree is $r$. This proves the assertion. For a line bundle on a scheme, applying the associated-bundle map gives the identical statement in its $K_0$. This does not assert that the two resulting bundles are nonisomorphic on every scheme.
\end{proof}

\begin{proposition}[G-04: An integral coefficient model and its exact specialization gate]\label{G-04}
Let $W_{\mathbb Z}$ be a finite free module, let $\rho:\mathrm{GL}_{N,\mathbb Z}\to\mathrm{GL}(W_{\mathbb Z})$ be a homogeneous polynomial representation of positive degree $e$, and let $f\in W_{\mathbb Z}$. Define $X_{\mathbb Z}$ as the affine schematic closure of its orbit. Its coordinate algebra is torsion-free, hence flat over $\mathbb Z$. In degree $d$ its orbit-coefficient map is an explicit finite integer matrix
\[
 A_d:\operatorname{Sym}^d(W_{\mathbb Z}^*)\longrightarrow
 \mathbb Z[g_{11},\ldots,g_{NN}]_{\leq ed},\qquad P\longmapsto P(\rho(g)f).
\]
Its rank over $\mathbb Q$ is the degree-$d$ coordinate-ring dimension. If $\delta_d$ is any nonzero maximal-rank minor of $A_d$, then this dimension and the degree-$d$ kernel commute with base change over $\mathbb Z[1/\delta_d]$. At a prime dividing $\delta_d$, rank must be tested anew; the fiber of the integral closure need not be the schematic orbit closure of the reduced point.
\end{proposition}
\begin{proof}
The coordinate homomorphism of the orbit map has target $\mathbb Z[\mathrm{GL}_N]=\mathbb Z[g_{ij},\det(g)^{-1}]$, a torsion-free $\mathbb Z$-module. Its kernel $I$ defines the schematic closure and its image $\mathbb Z[W]/I$ is torsion-free. A module over the principal ideal domain $\mathbb Z$ is flat if it is torsion-free, proving the first assertion. Since the representation is polynomial, each degree-$d$ homogeneous coordinate polynomial maps to a polynomial of degree at most $ed$ in the entries of $g$. Equality on invertible matrices is equality in the polynomial ring because localization by the nonzero polynomial $\det(g)$ is injective. This proves the finite coefficient-matrix description.
Let $r=\operatorname{rank}_{\mathbb Q}A_d$. All $(r+1)$-minors vanish as integers, while an $r$-minor $\delta_d$ is nonzero. Over $\mathbb Z[1/\delta_d]$, use that invertible block for row and column elimination. The vanishing of larger minors makes the remaining Schur complement zero, so the matrix is equivalent to $\operatorname{diag}(I_r,0)$. Its image and kernel are direct summands there and commute with every base change. Without inversion, an inclusion of lattices need not remain injective modulo a prime: already multiplication by $p$ on $\mathbb Z$ has rank one over $\mathbb Q$ and rank zero modulo $p$. Thus no unsupported all-primes orbit-fiber conclusion follows. Homogeneity sends degree $d$ to degree $ed$, so the kernel is homogeneous and the stated grading is well defined. If $r=0$, take $\delta_d=1$.
\end{proof}

\begin{proposition}[G-05: Ternary saturation and the three-row exclusion]\label{G-05}
Work over $\mathbb C$, and let $n\geq2$, $N=n^2$, and $\mathrm{Det}_n\subseteq\operatorname{Sym}^n(\mathbb C^N)^*$ be the determinant orbit closure. Then $\mathrm{Det}_n$ contains every degree-$n$ form supported on a fixed three-dimensional subspace. Consequently, for every partition $\lambda\vdash nd$ with $\ell(\lambda)\leq3$,
\[
 \operatorname{mult}_{S_\lambda\mathbb C^N}\mathbb C[\mathrm{Det}_n]_d
 =\operatorname{mult}_{S_\lambda\mathbb C^N}\operatorname{Sym}^d(\operatorname{Sym}^n\mathbb C^N).
\]
No multiplicity obstruction with at most three rows can separate any other orbit closure of a degree-$n$ form in the same ambient space, including a padded permanent, from $\mathrm{Det}_n$.
\end{proposition}
\begin{proof}
We use the established linear determinantal representation theorem for smooth plane curves: a smooth plane degree-$n$ curve together with a line bundle $L$ of degree $g-1$ with $H^0(L)=0$ admits an $n\times n$ matrix of linear forms whose determinant is its defining form up to a nonzero scalar (Beauville, 2000, Proposition 3.1(a)). Such an $L$ exists over $\mathbb C$: for $g\geq1$, the effective locus in $\operatorname{Pic}^{g-1}$ is the image of the symmetric power $C^{(g-1)}$, of dimension at most $g-1$, whereas $\operatorname{Pic}^{g-1}$ has dimension $g$; for $g=0$, take the degree-$-1$ line bundle. Scaling one row absorbs the nonzero scalar. Hence every smooth ternary degree-$n$ form is obtained by a linear substitution in $\det_n$. Any endomorphism substitution is a limit of invertible substitutions because $\mathrm{GL}_N$ is Zariski dense in $\operatorname{End}(\mathbb C^N)$; continuity of the polynomial action therefore puts these forms in $\mathrm{Det}_n$. Smooth ternary forms form a dense open set, so its closedness puts every ternary form there.
Let $W=\operatorname{Sym}^d(\operatorname{Sym}^n\mathbb C^N)$, and choose the Borel convention whose highest-weight vector has nonnegative weight $\lambda=(\lambda_1,\lambda_2,\lambda_3,0,\ldots,0)$. In a monomial basis of $W$, zero weight in positions $4,\ldots,N$ forces every coefficient variable to involve only the first three basis vectors. Thus the entire highest-weight space $W_\lambda^{\mathrm{hw}}$ is contained in $\operatorname{Sym}^d(\operatorname{Sym}^n\mathbb C^3)$. A nonzero polynomial in these coefficient variables cannot vanish on all ternary forms. By the first paragraph, the degree-$d$ determinant ideal therefore meets $W_\lambda^{\mathrm{hw}}$ trivially. Exactness and complete reducibility in characteristic zero give the asserted multiplicity equality. Every other closed subvariety has a degree-$d$ coordinate ring which is a quotient of $W$, so its multiplicity cannot exceed this common ambient value. This proves the exclusion. The assertion is about complex characteristic-zero multiplicities, not an unqualified statement in positive characteristic.
\end{proof}

\begin{proposition}[G-06: The four-variable gate is nonvacuous]\label{G-06}
For $n\geq4$, the closure $\mathcal D_{n,4}$ of degree-$n$ quaternary forms which are determinants of $n\times n$ matrices of linear forms has dimension at most $2n^2+2$. Since
\[
 \binom{n+3}{3}-(2n^2+2)=\binom{n-1}{3}>0,
\]
this is a proper subvariety of the space of quaternary forms. The determinant ideal in $\operatorname{Sym}(\operatorname{Sym}^n\mathbb C^{n^2})$ therefore has a nonzero isotypic component with a partition of exactly four rows in some finite degree.
\end{proposition}
\begin{proof}
Consider the polynomial map
\[
 (A_1,A_2,A_3,A_4)\longmapsto\det(x_1A_1+x_2A_2+x_3A_3+x_4A_4)
\]
from an affine space of dimension $4n^2$. The group
$H=\{(P,Q)\in\mathrm{GL}_n^2:\det(P)=\det(Q)\}$ acts by $A_i\mapsto PA_iQ^{-1}$, preserving this map. It has dimension $2n^2-1$. Its generic stabilizer has dimension one. Indeed, for a tuple with $A_1=I$, a diagonal $A_2$ with pairwise distinct eigenvalues, and an $A_3$ with every off-diagonal entry nonzero, a stabilizer has $Q=P$; commuting with $A_2$ makes $P$ diagonal, and commuting with $A_3$ makes it scalar. The scalar pairs $(cI,cI)$ always stabilize every tuple. Upper semicontinuity of stabilizer dimension, or the rank calculation for the commutation equations on this exhibited open set, gives generic dimension one. Thus generic fibers contain orbits of dimension $2n^2-2$, and the dimension-of-fibers theorem yields the stated image bound. The displayed identity follows by expanding the binomial coefficient.
Let $\pi$ set all but the first four variables to zero. For every form in $\mathrm{Det}_n$, its image under $\pi$ lies in $\mathcal D_{n,4}$: approximate the form by invertible substitutions of $\det_n$, apply $\pi$, and use the closedness of $\mathcal D_{n,4}$. Choose a nonzero homogeneous equation $P$ of this proper cone. Then $P\circ\pi$ is a nonzero homogeneous equation of $\mathrm{Det}_n$, using only first-four coefficient variables. The ideal of $\mathcal D_{n,4}$ is $\mathrm{GL}_4$-stable; choose a highest-weight vector in a nonzero irreducible submodule of its degree-$d$ part. Embedded in the full coordinate ring, it is also a $\mathrm{GL}_{n^2}$ highest-weight vector: raising operators involving a variable beyond the fourth annihilate it. Its highest weight has at most four rows, and G-05 rules out at most three. Hence its length is exactly four. Existence of such an equation does not establish that it is nonzero on the padded permanent.
\end{proof}

\begin{proposition}[G-07: An exact four-row determinant-closure formula]\label{G-07}
Set $D=320112$. For every $d\geq0$ and partition $\lambda\vdash4d$ with $\ell(\lambda)\leq4$, padded by zeros to length four, the degree-$d$ determinant-closure multiplicity satisfies
\[
 m_{\mathrm{Det}_4}(\lambda,d)
 =a_\lambda(d[4])-a_{\lambda-(D,D,D,D)}((d-D)[4]).
\]
The second term is defined to be zero if $d<D$ or $\lambda_4<D$. In particular the four-row sector is ambient for $d<D$, and at $d=D$ the multiplicity of $\lambda=(D^4)$ is exactly one less than its ambient plethysm multiplicity.
\end{proposition}
\begin{proof}
Leal, Lozano Huerta, and Vite (2024), Theorem 2 and Table 2, identify the closure of linear determinantal quartic surfaces as an irreducible reduced hypersurface of degree $D$ in $\mathbb P(\operatorname{Sym}^4\mathbb C^4)^*$. Table 2 identifies this component $F_1$ by $a=(5,5,5,5)$, $b=(6,6,6,6)$, so every matrix entry has degree one. Its affine cone has a prime principal ideal $(F)$ in the polynomial ring $A=\operatorname{Sym}(\operatorname{Sym}^4\mathbb C^4)$, with $F$ homogeneous of degree $D$.
The line $\mathbb CF$ is stable under $\mathrm{GL}_4$. Every algebraic one-dimensional character of $\mathrm{GL}_4$ is a power of the determinant. Since scalar $zI$ acts on $A_D$ by $z^{4D}$, this character is $\det^D$. Thus multiplication by $F$ gives the exact sequence of $\mathrm{GL}_4$-modules
\[
 0\longrightarrow \det^D\otimes A_{d-D}\xrightarrow{\ \cdot F\ }A_d
 \longrightarrow(A/(F))_d\longrightarrow0.
\]
For $d<D$ its first term is zero.
Let $\pi$ restrict forms in sixteen variables to their first four variables. The image $\pi(\mathrm{Det}_4)$ is precisely the affine cone just described. One inclusion follows by continuity from the orbit of $\det_4$. For the converse, every four-variable linear determinant is an endomorphism substitution in $\det_4$, hence belongs to its orbit closure; the same is true of every limit of such forms, and such a form is fixed by $\pi$. It follows that the intersection of the full determinant ideal with the first-four coefficient polynomial ring is exactly $(F)$. Highest-weight spaces for partitions of at most four rows use only these coefficient variables, by the weight argument of G-05. Applying complete reducibility to the exact sequence and shifting a highest weight by the determinant character proves the formula. At $d=D$ the source is $\det^D\otimes\mathbb C$, so the deficit is one.
\end{proof}
\paragraph{Source correction for G-07.}
The attached Landsberg source, section \texttt{zhanglee}, theorem \texttt{zlthm}, records $640224$. The later result used here is the degree of the reduced prime divisor. Leal--Lozano Huerta--Vite, section 4.1 immediately before Theorem 2, explicitly gives $P_{20,2}=\tfrac12D_{23,14}$. This is the precise cycle-multiplicity distinction; using the unreduced Noether--Lefschetz coefficient as the polynomial degree doubles the answer. This correction concerns the cited literature, not a purported visual correction of cote 133.

\begin{proposition}[G-08: A uniform low-variable determinantal construction]\label{G-08}
Let $m\geq1$ and $k\geq2$, put $a=\lfloor k/2\rfloor$, $b=\lceil k/2\rceil$, and
\[
 B(m,k)=\binom{m+a}{a}+\binom{m+b}{b}-1.
\]
Over any commutative ring, every degree-$m$ homogeneous polynomial $q$ in $k$ variables admits a matrix of linear forms of size $B(m,k)$ with determinant
\[
 \ell^{B(m,k)-m}q,
\]
where $\ell$ may first be an independent variable and then be specialized to any linear form. The construction is polynomial in the coefficients of $q$ and defined over $\mathbb Z$.
Consequently, over $\mathbb C$, whenever $n\geq B(m,k)$, no multiplicity obstruction of length at most $k$ can separate the orbit closure of an arbitrarily padded degree-$m$ form $\ell^{n-m}f$ from $\mathrm{Det}_n$.
\end{proposition}
\begin{proof}
Split the variables into disjoint lists $x=(x_1,\ldots,x_a)$ and $y=(y_1,\ldots,y_b)$, and write
\[
 q=\sum_{|\alpha|+|\beta|=m}c_{\alpha\beta}x^\alpha y^\beta.
\]
Make a rooted directed tree with one vertex for every $x$-monomial of degree at most $m$. The parent of a nonconstant monomial is obtained by deleting one occurrence of its largest-index variable; orient the edge away from the constant root and label it by that variable. There is exactly one root-to-vertex path, with label product $x^\alpha$. Make the corresponding $y$-tree with every edge reversed, so the unique path from vertex $\beta$ to its constant root has product $y^\beta$. Connect vertex $\alpha$ in the first tree to vertex $\beta$ in the second by an edge of constant label $c_{\alpha\beta}$ exactly when $|\alpha|+|\beta|=m$. Every source-to-sink path uses exactly one connecting edge; its label product is $c_{\alpha\beta}x^\alpha y^\beta$. Therefore its path polynomial is exactly $q$. The graph is acyclic and has
\[
 s=\binom{m+a}{a}+\binom{m+b}{b}=B(m,k)+1
\]
vertices.
Choose a topological ordering with source first and sink last, and let $A$ be its strictly upper triangular weighted adjacency matrix. Then $\det(I-A)=1$ and $(I-A)^{-1}=I+A+\cdots+A^{s-1}$. The $(1,s)$ entry is the source-to-sink path polynomial $q$. The adjugate identity therefore expresses $q$, up to the fixed sign $(-1)^{s+1}$, as the determinant of the $(s-1)\times(s-1)$ minor of $I-A$ obtained by deleting row $s$ and column $1$. Change the sign of one row if needed. Its entries are affine linear in $x,y$.
Replace each constant entry $c$ of this minor by $c\ell$. Its determinant is the degree-$B(m,k)$ homogenization of $q$, hence is $\ell^{B(m,k)-m}q$. This identity is universal over $\mathbb Z[c_{\alpha\beta},x,y,\ell]$. Padding the matrix with $n-B(m,k)$ diagonal entries $\ell$ proves the corresponding size-$n$ statement.
For the multiplicity assertion, restrict any invertible substitution of $\ell^{n-m}f$ to $k$ coordinate variables. The result has the form $L^{n-m}q$, with $q$ a homogeneous degree-$m$ form in these variables. The construction puts this result in $\mathrm{Det}_n$. Continuity gives the same inclusion for the entire padded orbit closure. If $P$ is any highest-weight polynomial of length at most $k$ vanishing on $\mathrm{Det}_n$, it uses only coefficient variables supported on the first $k$ underlying variables (in the sense made precise in G-05), and therefore vanishes on the other orbit closure as well. Thus the highest-weight space of the determinant ideal embeds in that of the padded-form ideal. Subtracting their dimensions from the common ambient highest-weight multiplicity gives $m_{\mathrm{padded}}(\lambda,d)\leq m_{\mathrm{Det}_n}(\lambda,d)$ in every degree.
\end{proof}
\paragraph{Quantified implication.}
For $k=4$, the construction gives $B(m,4)=(m+1)(m+2)-1$. More generally $B(m,k)=O_k(m^{\lceil k/2\rceil})$. Thus for a fixed number of rows this sector is excluded once the padding exceeds the displayed bound. This is an explicit restricted lower-bound obstruction theorem, not a general separation result and not an efficient solution of arbitrary plethysm.

\begin{proposition}[G-09: A sharper low-row exclusion by shared exponent chains]\label{G-09}
Let $m\geq1$ and $k\geq3$. Define the following integer bounds. If $k=2r+1$ is odd, set
\[
 C(m,k)=2\binom{m+r}{r}-1.
\]
If $k=2r+2\geq4$ is even, for $1\leq h\leq m+1$ put $Q=\lfloor m/h\rfloor$, $s_0=\min(h-1,m)$ and
\[
 C_h(m,k)=\sum_{q=0}^{Q}\binom{m-qh+r}{r}+Q(h-1)
          +\sum_{s=0}^{s_0}\binom{m-s+r}{r}-1,
 \qquad C(m,k)=\min_{1\leq h\leq m+1}C_h(m,k).
\]
Over a field, for every nonzero linear form $L$ and homogeneous degree-$m$ form $q$ in $k$ variables, $L^{C(m,k)-m}q$ has a linear determinantal representation of size $C(m,k)$. In coordinates with $L=x_0$ the construction is defined over $\mathbb Z$ and polynomial in the coefficients of $q$. For fixed $k\geq3$,
\[
 C(m,k)=O_k\bigl(m^{(k-1)/2}\bigr).
\]
Consequently, over $\mathbb C$, if $n\geq C(m,k)$, there is no multiplicity obstruction with at most $k$ rows separating any padded degree-$m$ form from $\mathrm{Det}_n$. This strengthens G-08 in the specified range.
\end{proposition}
\paragraph{Prior-result boundary.}
The asymptotic determinantal-representation order used here is established: Boralevi, van Doornmalen, Draisma, Hochstenbach, and Plestenjak (2017), Theorem 1.3, prove order $m^{v/2}$ for a fixed number $v\geq2$ of affine variables, with a matching generic lower bound. Their proof in section 4 uses connected monomial spaces and an even/odd binary-digit split. We do not claim that order as new. The proof below gives the displayed finite integer bound by a different shared-chain construction, records its coefficient-ring scope, and proves the resulting GCT multiplicity exclusion. These are the forms needed here that the cited theorem does not state.
\begin{proof}
Extend $L$ to a basis of linear forms and write it as $x_0$. Dehomogenize $q$ by setting $x_0=1$, obtaining an arbitrary polynomial $p$ of degree at most $m$ in $k-1$ variables. It suffices to construct an affine linear determinant of size $C(m,k)$ equal to $p$. Homogenizing each matrix entry by $x_0$ then gives $x_0^{C(m,k)-m}q$.

For odd $k=2r+1$, split the dehomogenized variables into lists $x$ and $y$ of $r$ variables each. Use the two monomial trees from G-08, each with $\binom{m+r}{r}$ vertices, and connect every pair $(\alpha,\beta)$ with $|\alpha|+|\beta|\leq m$ by its coefficient in $p$. Every source-to-sink path is unique for its two monomial indices, so its path polynomial is $p$. The cofactor construction in G-08 has size one less than the total number of vertices, giving the displayed bound.

For even $k=2r+2$, write the variables as $x_1,\ldots,x_r,y_1,\ldots,y_r,z$ and fix $h$. Decompose the exponent of $z$ uniquely as $e=qh+s$, $0\leq s<h$. On the left, make a directed chain with $Qh$ edges, all labelled $z$, beginning at the source. At its vertex of index $qh$ attach the outward $x$-monomial tree of degrees at most $m-qh$, identifying its constant root with that chain vertex. Every such tree vertex indexed by $(q,\alpha)$ has a unique source path, with product $z^{qh}x^\alpha$. Chain vertices whose indices are not multiples of $h$ are transit vertices and receive no connecting edges. The left graph has
\[
 (Qh+1)+\sum_{q=0}^Q\left[\binom{m-qh+r}{r}-1\right]
 =\sum_{q=0}^Q\binom{m-qh+r}{r}+Q(h-1)
\]
vertices.

On the right, make a chain with vertices $0,\ldots,s_0$, directed toward vertex $0$, the sink, and again label its edges by $z$. At vertex $s$ attach the inward $y$-monomial tree of degrees at most $m-s$, identifying the root with that vertex. Every tree vertex indexed by $(s,\beta)$ has a unique sink path with product $y^\beta z^s$. The right graph has $\sum_{s=0}^{s_0}\binom{m-s+r}{r}$ vertices. Connect left vertex $(q,\alpha)$ to right vertex $(s,\beta)$ with the constant coefficient of $x^\alpha y^\beta z^{qh+s}$ in $p$ when $|\alpha|+|\beta|+qh+s\leq m$.

All connecting edges go from left to right. Thus every source-to-sink path uses exactly one such edge, and the uniqueness of the two tree paths and the division $e=qh+s$ gives exactly the required monomial with its coefficient, once. The path polynomial is $p$. Taking the same cofactor of $I-A$ as in G-08 yields an affine linear determinant of size $C_h(m,k)$. This graph has a path capable of realizing $z^m$ and a connecting edge, so its number of vertices is at least $m+2$; hence $C_h(m,k)\geq m+1$. The odd bound is also greater than $m$. Homogenization therefore uses a nonnegative exponent. Minimizing over the finite choices of $h$ proves the stated representation bound.

For the asymptotic statement in the even case take $h=\lceil\sqrt m\rceil$. The number of binomial summands is at most $m/h+1+h\leq2\sqrt m+2$, each bounded by $\binom{m+r}{r}=O_r(m^r)$, while $Q(h-1)\leq m$. Thus $C(m,k)=O_r(m^{r+1/2})$. The odd case is immediate from its binomial formula. These exponents are $(k-1)/2$ in both cases.

For the GCT consequence restrict a substituted padded form to any fixed $k$ underlying variables. If its restricted padding $L$ is nonzero, apply the determinant construction and pad with $n-C(m,k)$ diagonal copies of $L$. If $L=0$, the form is zero because $n\geq C(m,k)>m$. Thus every restriction lies in $\mathrm{Det}_n$. Apply the highest-weight ideal argument from G-08, which uses exactly this containment, to obtain the multiplicity inequality in every degree.
\end{proof}
\paragraph{Integral scope and growth of the remaining gate.}
With $L=x_0$ the graph and cofactor identity are universal over the coefficient ring. Extending an arbitrary nonzero $L$ to coordinates is a field operation; over a general base one instead works on a chart where a coefficient of $L$ is invertible and chooses the resulting local basis. No global polynomial choice of that basis is asserted. If $n=m^c$ with fixed $c>0$, every fixed $k\geq3$ satisfying $(k-1)/2<c$ is excluded for sufficiently large $m$. This is a quantitative restriction on possible multiplicity obstructions, not a claim that the remaining higher-row obstruction exists.
